% 对应题号 M2-015，本题条件、变量及几何关系逐项核对。
% 修订：启用中文断行；核对题设及图文对应。
\documentclass[tikz,border=14pt]{standalone}
\usepackage{fontspec,xeCJK}
\usepackage{amsmath,amssymb}
\usepackage{xcolor}
\usepackage{tikz}
\usetikzlibrary{arrows.meta,calc,patterns,positioning,bending}
\IfFontExistsTF{Heiti SC}{\IfFontExistsTF{Heiti SC}{\setmainfont{Heiti SC}}{\setmainfont{FandolHei-Regular}}\IfFontExistsTF{Heiti SC}{\setsansfont{Heiti SC}}{\setsansfont{TeX Gyre Heros}}}{\IfFontExistsTF{Songti SC}{\setmainfont{Songti SC}}{\setmainfont{FandolSong-Regular}}\IfFontExistsTF{Songti SC}{\setsansfont{Songti SC}}{\setsansfont{FandolSong-Regular}}}
\definecolor{bluefield}{RGB}{24,76,255}
\definecolor{redfield}{RGB}{232,49,35}
\definecolor{orangefield}{RGB}{236,132,18}
\definecolor{greenfield}{RGB}{18,132,86}
\definecolor{purplefield}{RGB}{112,72,190}
\definecolor{linegray}{RGB}{60,68,78}
\definecolor{softblue}{RGB}{238,243,255}
\definecolor{softgreen}{RGB}{238,249,244}
\definecolor{softorange}{RGB}{255,247,235}
\tikzset{
  arrow/.style={-{Latex[length=3.1mm,width=2.1mm]},line width=1.1pt},
  thinarr/.style={-{Latex[length=2.5mm,width=1.7mm]},line width=.85pt},
  axis/.style={-{Latex[length=2.5mm,width=1.7mm]},line width=.8pt,draw=linegray},
  guide/.style={draw=linegray!55,dashed,line width=.65pt},
  box/.style={draw=linegray,rounded corners=3pt,fill=white,align=center,inner sep=6pt,line width=.75pt},
  bluebox/.style={box,draw=bluefield,fill=softblue},
  greenbox/.style={box,draw=greenfield,fill=softgreen},
  orangebox/.style={box,draw=orangefield,fill=softorange}
}
\IfFontExistsTF{Heiti SC}{\setCJKmainfont[AutoFakeBold=1.5]{Heiti SC}}{\setCJKmainfont[AutoFakeBold=1.5]{FandolHei-Regular}}
\IfFontExistsTF{Heiti SC}{\setCJKsansfont{Heiti SC}}{\setCJKsansfont{FandolHei-Regular}}
\begin{document}

\begin{tikzpicture}[x=1cm,y=1cm]
\path[use as bounding box] (-10,-4.6) rectangle (10,4.6);
\fill[white] (-10,-4.6) rectangle (10,4.6);
\node[font=\Large\bfseries,text width=18.6cm,align=center] at (0,4.05) {异号平行线电荷：横截面中等势圆与场线圆正交};
\begin{scope}[shift={(-4.8,-.15)}]
\begin{scope}\clip(-3,-2.8)rectangle(3,2.8);
\foreach\s in{.5,1,1.7}{\pgfmathsetmacro{\c}{sqrt(1.4^2+\s^2)}\draw[bluefield](\c,0)circle(\s);\draw[redfield](-\c,0)circle(\s);}
\foreach\b in{-1.6,-.6,0,.6,1.6}{\pgfmathsetmacro{\r}{sqrt(1.4^2+\b*\b)}\draw[greenfield!65](0,\b)circle(\r);}
\end{scope}
\fill[bluefield](1.4,0)circle(.08)node[above right,fill=white]{$-\lambda$};\fill[redfield](-1.4,0)circle(.08)node[above left,fill=white]{$+\lambda$};
\node[font=\small,fill=white,inner sep=2pt]at(0,3.1){红、蓝：等势圆；绿：场线圆};
\end{scope}
\begin{scope}[shift={(4.65,-.1)}]
\node[bluebox,text width=6.8cm] at (0,1.65) {$\displaystyle \varphi=\frac{\lambda}{2\pi\varepsilon_0}\ln\frac{r_-}{r_+}$};
\node[greenbox,text width=6.8cm] at (0,0) {等势线满足 $r_-/r_+=\text{常量}$；场线与之处处正交。};
\node[orangebox,text width=6.8cm] at (0,-1.7) {这是三维平行线电荷的二维横截面，不是两个点电荷模型。};
\end{scope}
\end{tikzpicture}
\end{document}
